% !TeX root = ../../main.tex
% chapters/4_courb/intro.tex -- one section of Chapter 4 (CoUrb 2026, published).
% Split out of chapters/4_courb.tex on 2026-07-28 (author request: the paper
% chapters were single long files, while the originals under articles/CoUrb_2026/src_en/sections/
% are one file per section). The split is PURELY MECHANICAL: content was cut at
% \section boundaries and not edited. The chapter file now \inputs these in order.
% Editing rules are unchanged and still governed by the chapter's status above:
% see NORTH_STAR section 4 before changing any sentence here.
\section{Introduction}
\label{sec:courb:intro}

Location-Based Social Networks (LBSNs), such as Gowalla \cite{cho2011gowalla,jure2014snap}, record user check-ins with geographic location, visit time, and category of the visited place. In this context, a Point of Interest (POI) corresponds to a semantically identifiable physical place, such as restaurants, stores, airports, or parks. The analysis of these data is relevant to applications such as recommendation and urban mobility. In this scenario, two tasks stand out:

\begin{enumerate}
    \item \textbf{POI Category Classification}: classify the semantic category (e.g., \textit{Food}, \textit{Shopping}, \textit{Travel}) of a POI from its features.
    \item \textbf{Next-POI Prediction}: predict the category of the next POI that a user will visit, given the historical sequence of their check-ins.
\end{enumerate}

In principle, these tasks may be treated as related, since they share the same mobility dataset and the same semantic space of categories, which makes multi-task learning (MTL) a promising strategy to exploit information shared between them \cite{caruana1997multitask}. However, they impose distinct demands. POI Category Classification is a non-sequential task, since it depends mainly on the attributes of the POI itself and on its spatial context. Next-POI Prediction is explicitly sequential, requiring the modeling of temporal order, recurrence, and transitions along the user trajectory. Thus, the information useful for one task may not benefit the other in the same way.

MTLnet, proposed in \cite{silva2025mtlnet}, exploits this potential by modeling the two tasks in a single multi-task architecture with hard parameter sharing. However, the model uses exclusively vector representations based on the graph structure, generated by \textit{Deep Graph Infomax} (DGI), which encodes spatial and categorical information in a single 64-dimensional vector, without explicitly incorporating continuous geographic coordinates or temporal visitation patterns.

Given this, the question arises of whether decomposing the input into specialized encoders for space, time, and category produces a base more suitable for the two tasks than the monolithic embedding of DGI. The hypothesis is that decoupled representations capture complementary dimensions of human mobility that are not explicitly modeled by the original approach.

In this sense, this chapter proposes ST-MTLNet (\textit{Spatial-Temporal MTLNet}), which integrates three independent representations: a continuous spatial representation for geographic coordinates, a temporal representation (Time2Vec) for visitation patterns, and a hierarchical categorical representation for regional and structural relationships between POIs. In the spatial dimension, we evaluate two architectures with distinct assumptions, SIREN and Sphere2Vec-M, originally validated in geospatial tasks of remote sensing and ecology \cite{wu2024torchspatial}, but still little explored in multi-task learning of POIs in LBSNs.

The experiments were carried out with the public Gowalla dataset \cite{jure2014snap} in the states of Florida, California, and Texas. In categorical classification, ST-MTLNet outperforms MTLnet in all evaluated category-state combinations, with average gains per state of 20.2 to 22.0 percentage points, considering the better of the two spatial encoders in each combination. In Next-POI Prediction, the proposed model outperforms the baseline in most scenarios, with emphasis on categories such as \textit{Food}. The comparison between SIREN and Sphere2Vec-M also shows that, although the modular approach is consistently superior to the baseline, the most suitable spatial encoder depends on the geographic distribution of the POIs in each territory.

Therefore, the main contributions of this chapter are:
\begin{itemize}
\item The proposal of ST-MTLNet, which incorporates continuous spatial, temporal, and hierarchical categorical representations into MTLnet.

\item The comparative evaluation of SIREN and Sphere2Vec-M in three states with distinct territorial structures.

\item The public release of the source code and the experimental pipeline.\footnote{The complete source code is available at \url{https://github.com/TarikSalles/Spatial_Embeddings}}
\end{itemize}

This chapter is organized as follows: Section \ref{sec:courb:related} presents the related work. Section \ref{sec:courb:methods} describes the proposed methodology. Section \ref{sec:courb:results} discusses the experimental results. Finally, Section \ref{sec:courb:conclusion} presents the conclusions and future work.
